\documentclass[10pt,a4paper]{amsart}
\usepackage[margin=19mm]{geometry}
\usepackage{amsmath,amssymb,mathtools}
\usepackage[scr=rsfs,cal=euler]{mathalfa}
\usepackage{xfrac}
\usepackage[unicode,hidelinks,bookmarks=false]{hyperref}
\newcommand{\ie}{i.e.\ }
\newcommand{\cf}{cf.\ }
\newcommand{\resp}{respectively\ }
\newcommand{\Subsection}{\subsection}
\providecommand{\nobreakdash}{\nobreak\hbox{-}}
\setlength{\emergencystretch}{2em}
\multlinegap0pt
\usepackage{bm}
\usepackage{bbm}
\usepackage{dsfont}
\usepackage{xspace}
\usepackage{cancel}
\usepackage{mathtools}
\usepackage{amsthm}
\makeatletter
\@ifundefined{theorem}{\newtheorem{theorem}{Theorem}}{}
\makeatother
\title[Set Valued Riemann-Liouville integral and some Regular Selec]{Set Valued Riemann-Liouville integral and some Regular Selections}
\author{Subhash Chandra, Syed Abbas}
\date{Source: arXiv:2512.23584v1 (CC BY 4.0)}
\begin{document}
\maketitle

\noindent\emph{Source and licence.} Set Valued Riemann-Liouville integral and some Regular Selections. Version arXiv:2512.23584v1 (CC BY 4.0), distributed under the Creative Commons Attribution 4.0 International licence, as recorded in the arXiv licence metadata for this exact version and shown on the abstract page https://arxiv.org/abs/2512.23584v1. Full rights notice: licenses/arxiv-2512.23584-CC-BY-4.0.txt.\par\smallskip

\noindent\textit{Editorial scope.} Counted unit: the Theorem whose source label is \texttt{conti} in the article (source0.tex lines 343--348, source label \texttt{conti}), delivered together with the article's own complete original proof of it (source0.tex lines 349--471, one \texttt{proof} environment of 123 lines). No mathematics is written, repaired or re-derived: statement and proof are the article's own text line for line. Required context retained uncounted: THnonempty (lines 298--301). Every definition, display and label of those lines is retained unchanged. Omitted from this package and not redistributed: 1--297, 302--342, 472--852. Nothing inside a retained excerpt is omitted or altered: every retained body is delivered exactly as the article's lines are written, so no omission receipt of any kind accompanies this package. Citations: the retained excerpts carry no citation command at all, so no bibliography entry of the article is transcribed and no reference is redistributed. Reference adaptation: none was needed and none was made; every reference inside the delivered unit resolves inside it, so no unresolved reference or citation is produced. Printed slips kept literal: no printed word, symbol, hypothesis or number of the counted statement or of its proof was altered. Layout and class: the article's own class is \texttt{amsart} with packages including amsthm, mathrsfs, amssymb, amsmath, enumerate, mathtools, breqn, hyperref, graphicx, enumitem, xcolor; the wrapper is the lane's pinned \texttt{amsart} editorial wrapper at 10pt on A4, which restates the theorem-like environments the retained text needs and adds the article's own macro definitions that the retained text needs (none), with extra wrapper packages (bm, bbm, dsfont, xspace, cancel, mathtools). The delivered numbering is the unit's own. Assets: no retained span contains a figure environment, a graphics-inclusion command, a PDF-page inclusion, a table or a TikZ picture; no image file is copied into this package, the package declares no asset dependency and no image file is redistributed with it. Licence: version arXiv:2512.23584v1 of this article is distributed under the Creative Commons Attribution 4.0 International licence, as recorded in the arXiv licence metadata for this exact version and shown on the abstract page https://arxiv.org/abs/2512.23584v1. A full rights notice is delivered at licenses/arxiv-2512.23584-CC-BY-4.0.txt. Characters: every retained excerpt is pure ASCII, so the delivered engine, XeLaTeX with its default Latin Modern text font, reports no missing character.
\begin{theorem}\label{THnonempty}
Let $F:[a,b]\rightrightarrows\mathbb{R}$ be a set-valued mapping with nonempty values. Assume that $F$ is Borel-measurable and integrably
bounded. Then for every $\rho>0$ and every $u\in[a,b]$, the set ${}_{a}\mathfrak{J}^{\rho}F(u)$ is nonempty.
\end{theorem}

\begin{theorem}\label{conti}
   Let $F:[a,b]\rightrightarrows\mathbb{R}$ be a set-valued mapping with nonempty compact values, Borel measurable and integrably bounded on $[a,b]$.
Let $\rho>0$. Then, for each $u\in[a,b]$, the set ${}_{a}\mathfrak{J}^{\rho}F(u)$ is nonempty and compact,  and the mapping
$$u \longmapsto {}_{a}\mathfrak{J}^{\rho}F(u)$$
 is continuous on $[a,b]$ with respect to $H_d$.
\end{theorem}

\begin{proof}
Since $F$ is Borel-measurable, integrably bounded and has nonempty compact values, Theorem \ref{THnonempty} guarantees that, for each $u\in[a,b]$, the set of
integrable selections of $F$ is nonempty and therefore
${}_{a}\mathfrak{J}^{\rho}F(u)\neq\emptyset$.
Let $\mathfrak{F}$ denote the family of all integrable selections of $F$. By integrable boundedness, there exists $h\in L^{1}([a,b])$ such that
$$\sup\{|x|:x\in F(t)\}\le h(t)
\quad\text{for a.e. }t\in[a,b].$$
In particular, for each $f\in\mathfrak{F}$ we have
\begin{equation}\label{EQ1}
|f(t)|\le h(t)\quad\text{for a.e. }t\in[a,b].
\end{equation}
For $f\in\mathfrak{F}$ and $u\in[a,b]$, write
$$I_f(u):=\frac{1}{\Gamma(\rho)}\int_{a}^{u}(u-t)^{\rho-1}f(t)dt,$$
so that
$${}_{a}\mathfrak{J}^{\rho}F(u)=\{I_f(u):f\in\mathfrak{F}\}.$$
For any $f\in\mathfrak{F}$ and $u\in[a,b]$, by using \eqref{EQ1} we have
\begin{align*}
    |I_f(u)|
\le
\frac{1}{\Gamma(\rho)}\int_{a}^{u}(u-t)^{\rho-1}|f(t)|dt
\le
\frac{1}{\Gamma(\rho)}\int_{a}^{u}(u-t)^{\rho-1}h(t)dt.
\end{align*}

Since $\rho>0$ and
$h\in L^{1}([a,b])$, the right-hand side is finite. 
Thus there exists $M>0$ such that

\begin{align*}
    |I_f(u)|\le M
\quad\text{for all }u\in[a,b]\text{ and all }f\in\mathfrak{F},
\end{align*}
and therefore ${}_{a}\mathfrak{J}^{\rho}F(u)$ is bounded for each $u$. To see that ${}_{a}\mathfrak{J}^{\rho}F(u)$ is closed, let  $\{I_{f_n}(u)\}_{n\in\mathbb{N}}\subset {}_{a}\mathfrak{J}^{\rho}F(u)$  be a sequence converging to some $y\in\mathbb{R}$.   Since each $f_n$ is an integrable selection of $F$ and $F$ is integrably
bounded, there exists a function $h\in L^{1}([a,b])$ such that
\begin{equation}\label{EQ2}
|f_n(t)|\le h(t)
\qquad\text{for a.e. }t\in[a,b]\text{ and for all }n\in\mathbb{N}.
\end{equation}
Consequently, by passing to a subsequence if necessary, we may assume that
\[
f_n(t)\longrightarrow f(t)
\quad\text{for a.e. }t\in[a,b],
\]
where $f\in L^{1}([a,b])$. Since $f_n(t)\in F(t)$ for a.e.~$t$ and each
$F(t)$ is compact,hence closed, it follows that $f(t)\in F(t)$ for a.e.~$t\in[a,b]$. Therefore, $f$ is also a selection of $F$, that is,
$f\in\mathfrak{F}$.

Now, by \eqref{EQ2} and the fact that the kernel 
$(u-t)^{\rho-1}/\Gamma(\rho)$ is nonnegative and belongs to $L^{1}([a,u])$,
the Dominated Convergence Theorem yields
\[
I_{f_n}(u)
=
\frac{1}{\Gamma(\rho)}\int_{a}^{u}(u-t)^{\rho-1}f_n(t)dt
\;\longrightarrow\;
\frac{1}{\Gamma(\rho)}\int_{a}^{u}(u-t)^{\rho-1}f(t)dt
=
I_f(u).
\]
Since $I_{f_n}(u)\to y$, it follows that $y=I_f(u)$, and hence  $y\in{}_{a}\mathfrak{J}^{\rho}F(u)$.  
Therefore ${}_{a}\mathfrak{J}^{\rho}F(u)$ is closed.  
Being closed and bounded in $\mathbb{R}$, it is compact.  
Together with nonemptiness from Theorem \ref{THnonempty}, this shows that  ${}_{a}\mathfrak{J}^{\rho}F(u)$ is a nonempty compact set for every  $u\in[a,b]$.
Let $u,v\in[a,b]$ and $f\in\mathfrak{F}$ be arbitrary. Without loss of generality assume $u\le v$. Then
\[
\begin{aligned}
I_f(v)-I_f(u)
&=
\frac{1}{\Gamma(\rho)}
\Bigg(
\int_{a}^{v}(v-t)^{\rho-1}f(t)dt
-
\int_{a}^{u}(u-t)^{\rho-1}f(t)dt
\Bigg) \\
&=
\frac{1}{\Gamma(\rho)}
\Bigg(
\int_{a}^{u}\big[(v-t)^{\rho-1}-(u-t)^{\rho-1}\big]f(t)dt
+
\int_{u}^{v}(v-t)^{\rho-1}f(t)dt
\Bigg).
\end{aligned}
\]
From \eqref{EQ1}, we have
$$\big|I_f(v)-I_f(u)\big| \le \Phi(u,v),$$
where
\[
\Phi(u,v) := \frac{1}{\Gamma(\rho)}
\Bigg( \int_{a}^{u}\big|(v-t)^{\rho-1}-(u-t)^{\rho-1}\big|h(t)dt + \int_{u}^{v}(v-t)^{\rho-1}h(t)dt \Bigg).
\]
We now show that $\Phi(u,v)\to 0$ as $v\to u$. For the first integral,
note that for fixed $u\in[a,b]$ and for a.e.~$t\in[a,u)$,
\[ (v-t)^{\rho-1}\to (u-t)^{\rho-1}
\quad\text{as }v\to u, \]
and hence
\[ \big|(v-t)^{\rho-1}-(u-t)^{\rho-1}\big|
\to 0 \quad\text{for a.e. }t\in[a,u]. \] 
 Moreover, since $\rho>0$, we have
\[ \big|(v-t)^{\rho-1}-(u-t)^{\rho-1}\big|
\le (v-t)^{\rho-1}+(u-t)^{\rho-1}, \]
and both $(v-t)^{\rho-1}$ and $(u-t)^{\rho-1}$ belong to $L^{1}([a,u])$ as functions of $t$. Thus
\[ \big|(v-t)^{\rho-1}-(u-t)^{\rho-1}\big|h(t) \le\big[(v-t)^{\rho-1}+(u-t)^{\rho-1}\big]h(t)\in L^{1}([a,u]), \]
and the Dominated Convergence Theorem yields
\[ \int_{a}^{u}\big|(v-t)^{\rho-1}-(u-t)^{\rho-1}\big|\,h(t)dt \longrightarrow 0 \quad\text{as }v\to u. \]

For the second integral, observe that for $v>u$,
\[ 0\le\int_{u}^{v}(v-t)^{\rho-1}h(t)dt \le\int_{u}^{v}(v-t)^{\rho-1}h(t)dt, \] and the right-hand side tends to $0$ as $v\to u$ by the absolute continuity of the Lebesgue integral, since $h\in L^{1}([a,b])$ and the interval $(u,v)$ shrinks to a point. Consequently,
\[ \Phi(u,v)\to 0 \quad\text{whenever }v\to u. \]
Now, fix $u,v\in[a,b]$ and take any $y\in {}_{a}\mathfrak{J}^{\rho}F(u)$. Then
$y=I_f(u)$ for some $f\in\mathfrak{F}$. Let $z:=I_f(v)\in
{}_{a}\mathfrak{J}^{\rho}F(v)$. By the above,
\[ |y-z|=\big|I_f(u)-I_f(v)\big|\le\Phi(u,v), \]
and hence
\[\sup_{y\in {}_{a}\mathfrak{J}^{\rho}F(u)} \inf_{z\in {}_{a}\mathfrak{J}^{\rho}F(v)}|y-z| \le\Phi(u,v). \]
Exchanging the roles of $u$ and $v$ yields
\[ \sup_{z\in {}_{a}\mathfrak{J}^{\rho}F(v)} \inf_{y\in {}_{a}\mathfrak{J}^{\rho}F(u)}|z-y| \le\Phi(u,v). \]
By the definition of the Hausdorff distance, we therefore have
\[ H_d\big({}_{a}\mathfrak{J}^{\rho}F(u),{}_{a}\mathfrak{J}^{\rho}F(v)\big) \le\Phi(u,v)
\quad\text{for all }u,v\in[a,b]. \]
Since $\Phi(u,v)\to 0$ as $v\to u$, this shows that the mapping
\[u\longmapsto {}_{a}\mathfrak{J}^{\rho}F(u) \]
is continuous on $[a,b]$ with respect to the Hausdorff distance $H_d$. This completes the proof. 
\end{proof}

\end{document}
